\begin{abstract}

Exploratory data analytics (EDA) is a sequential decision making process where analysts choose subsequent queries that might lead to some interesting insights based on the previous queries and corresponding results.
Data processing systems often execute the queries on samples to produce results with low latency. 
%To make interactions faster, several recent papers proposed use of samples to reduce the query latency. %
%However sampling creates approximation errors and can mislead the users in an interactive data exploration flow. This is because the results of the previous query can get distorted due to sampling and may prompt the user to a false path of analysis, which would have been useless if the original data was used. 
%There are several sampling techniques available. 
Different downsampling strategy preserves different statistics of the data and have different magnitude of latency reductions. The optimum choice of sampling strategy often depends on the particular context of the analysis flow and the hidden intent of the analyst. 
%
In this paper, we are the first to consider the impact of sampling in interactive data exploration settings as they introduce approximation errors. We propose a Deep Reinforcement Learning (DRL) based framework which can optimize the sample selection in order to keep the analysis and insight generation flow intact. 
%To model the user behavior with implicit intents in EDA, we further propose a novel intent-aware reward function to achieve more sample-efficient DRL.
Evaluations with 3 real datasets show that our technique can preserve the original insight generation flow while improving the interaction latency, compared to baseline methods.

\end{abstract}